% ============================================================
%  CHAPTER 5: Bulletproofs
% ============================================================

\chapter{Bulletproofs}
\label{chap:bulletproofs}

Bulletproofs were introduced by Bunz, Bootle, Boneh, Poelstra,
Wuille, and Maxwell in 2018 \cite{bulletproofs2018}. The defining
feature of Bulletproofs is that they require no trusted setup at
all. The public parameters are just $n$ random curve points that
anyone can generate from a public source of randomness, with no
secret toxic waste involved.
Unlike KZG and MKZG, Bulletproofs do not use pairings. Instead they
use a recursive inner product argument that halves the problem size
at each round. The result is a proof of $O(\log n)$ group elements,
but verification requires $O(n)$ scalar multiplications. This makes
Bulletproofs slower to verify than pairing-based schemes at large
sizes, but eliminates the trusted setup entirely.

\section{Bulletproofs Construction}

\textbf{Commitment:}

Given a scalar vector $\mathbf{u} \in \mathbb{F}^n$ and public
generators $\mathbf{G} \in \mathbb{G}^n$, the commitment is:

\begin{equation}
    C = \langle \mathbf{u}, \mathbf{G} \rangle
      = \sum_{i=0}^{n-1} u_i \cdot G_i
\end{equation}
\\
\textbf{Recursive Halving:}

The proof is built over $\log_2 n$ rounds. At each round the prover
splits the vectors in half and computes two cross terms:

\begin{equation}
    V_L = \langle \mathbf{u}_L, \mathbf{G}_R \rangle, \qquad
    V_R = \langle \mathbf{u}_R, \mathbf{G}_L \rangle
\end{equation}

The prover sends $V_L, V_R$ and receives challenge $\alpha$. The
vectors are then folded:

\begin{align}
    \mathbf{u}' &= \alpha \cdot \mathbf{u}_L + \alpha^{-1} \cdot \mathbf{u}_R \\
    \mathbf{G}' &= \alpha^{-1} \cdot \mathbf{G}_L + \alpha \cdot \mathbf{G}_R
\end{align}

The invariant $C' = \langle \mathbf{u}', \mathbf{G}' \rangle$ is
preserved. After $\log_2 n$ rounds, the vectors reduce to a single
scalar $\bar{u}$ and point $\bar{G}$, which are sent as the final
proof elements.
\\
\textbf{Verification:}

The verifier replays the folding using the round messages. At each
round the commitment is updated as:
\begin{equation}
    C \gets C + \alpha^2 \cdot V_L + \alpha^{-2} \cdot V_R
\end{equation}

and the generator vector is folded the same way. At the end the
verifier checks:

\begin{equation}
    C \stackrel{?}{=} \bar{u} \cdot \bar{G}
    \quad \text{and} \quad
    \bar{G} \stackrel{?}{=} G_{\mathsf{final}}
\end{equation}

Because the verifier must refold the full $\mathbf{G}$ vector at
each round, verification costs $O(n)$ scalar multiplications.
\\
\textbf{Proof Size:}

Each round contributes two curve points. The full proof is:
\begin{equation}
    \text{proof size}
    = \log_2(n) \times 2 \times |\mathbb{G}|
    + |\mathbb{F}| + |\mathbb{G}|
\end{equation}

\section{Pseudocode}

\begin{algorithm}[H]
\caption{Bulletproofs Setup, Commit, Prove, Verify}
\label{alg:bp-all}
\begin{algorithmic}[1]
\Statex \textbf{Setup}$(n)$: \Return $n$ random points $\mathbf{G}$
    \Comment{No secret involved}
\Statex \textbf{Commit}$(\mathbf{G}, \mathbf{u})$:
    \Return $\mathsf{MSM}(\mathbf{G}, \mathbf{u})$
\Statex
\Statex \textbf{Prove}$(C, \mathbf{u}, \mathbf{G})$:
\State $\mathsf{tr}.\mathsf{append}(C)$;\; $\mathsf{rounds} \gets []$
\While{$|\mathbf{u}| > 1$}
    \State Split $\mathbf{u}, \mathbf{G}$ into left/right halves
    \State $V_L \gets \mathsf{MSM}(\mathbf{G}_R, \mathbf{u}_L)$;\;
           $V_R \gets \mathsf{MSM}(\mathbf{G}_L, \mathbf{u}_R)$
    \State $\mathsf{tr}.\mathsf{append}(V_L, V_R)$;\;
           $\alpha \gets \mathsf{tr}.\mathsf{challenge}()$
           \Comment{Fiat-Shamir via SHA-256}
    \State $\mathbf{u} \gets \alpha\mathbf{u}_L + \alpha^{-1}\mathbf{u}_R$;\;
           $\mathbf{G} \gets \alpha^{-1}\mathbf{G}_L + \alpha\mathbf{G}_R$
    \State $\mathsf{rounds}.\mathsf{push}(V_L, V_R)$
\EndWhile
\State \Return $(\mathsf{rounds}, \mathbf{u}[0], \mathbf{G}[0])$
\Statex
\Statex \textbf{Verify}$(\mathbf{G}, C, \mathsf{rounds}, \bar{u}, \bar{G})$:
\State $\mathsf{tr}.\mathsf{append}(C)$
\For{\textbf{each} $(V_L, V_R)$ in $\mathsf{rounds}$}
    \State $\mathsf{tr}.\mathsf{append}(V_L, V_R)$;\;
           $\alpha \gets \mathsf{tr}.\mathsf{challenge}()$
    \State $C \gets C + \alpha^2 V_L + \alpha^{-2} V_R$;\;
           $\mathbf{G} \gets \alpha^{-1}\mathbf{G}_L + \alpha\mathbf{G}_R$
\EndFor
\State \Return $C = \bar{u} \cdot \bar{G}$ \textbf{and} $\mathbf{G}[0] = \bar{G}$
\end{algorithmic}
\end{algorithm}

\section{Implementation and Optimisations}

The implementation is in Rust using the \texttt{arkworks} library
\cite{arkworks}. Since Bulletproofs use no pairings, the generic
parameter is \texttt{CurveGroup} rather than \texttt{Pairing}, and
the same code runs on all three curves.

Two significant optimisations are present in the implementation that
differ from the textbook description of the protocol.

\subsection{Fiat-Shamir via SHA-256}

The original Bulletproofs theory describes an interactive protocol where the verifier
sends a fresh random $\alpha$ at each round. The implementation
replaces this with the Fiat-Shamir heuristic using SHA-256:

\begin{itemize}
    \item A \texttt{Transcript} struct accumulates all messages
          seen so far as byte vectors.
    \item At each round, after appending $V_L$ and $V_R$, the
          challenge is derived as:
          \texttt{alpha = SHA-256(all transcript bytes)},
          interpreted as a field element via
          \texttt{F::from\_le\_bytes\_mod\_order}.
    \item The verifier replays the exact same hash computation
          using the proof's round messages, producing the same
          challenges without any interaction.
\end{itemize}

This makes the proof a single transferable object. The prover runs
once and the output can be verified by anyone, anytime, without
any back-and-forth communication. This is the standard approach in
all practical deployments of Bulletproofs.

\subsection{Pippenger's MSM for Inner Products}

Instead of treating $\langle \mathbf{u}, \mathbf{G} \rangle$ as a
naive sum of $n$ scalar multiplications, costing $O(n)$ group
operations, the implementation uses
\texttt{VariableBaseMSM::msm\_bigint} from arkworks, which runs
Pippenger's algorithm:

\begin{itemize}
    \item Scalars are converted to \texttt{BigInt} representation
          before the call, which is required by the arkworks MSM
          interface.
    \item Points are converted to affine coordinates once before
          the MSM and reused across rounds where possible.
    \item The effective cost is $O(n / \log n)$ group operations
          rather than $O(n)$, giving a meaningful speedup for
          large $n$.
\end{itemize}

This optimisation applies in three places: the initial commit,
the $V_L$ computation at each round, and the $V_R$ computation
at each round. The total number of MSM calls across all rounds
is $2 \log_2 n$, with sizes $n, n/2, n/4, \ldots, 2$.

\subsection{Proof Size Calculation}

Proof size is computed as
$\log_2(n) \times 2 \times |\mathbb{G}.\mathsf{Affine}|
+ |\mathbb{G}.\mathsf{ScalarField}|
+ |\mathbb{G}.\mathsf{Affine}|$
using \texttt{std::mem::size\_of}. For BLS12-381 and BLS12-377
each affine point is 96 bytes and each scalar is 32 bytes. For
BN-254 each affine point is 64 bytes.
\\\\
\textbf{Benchmark Setup:}

Benchmarks are averaged over 10 runs. Sizes tested:
$n \in \{4, 8, 16, 32, 64, 128, 256, 512, 1024\}$.
MSM time is benchmarked separately at each size since it is the
dominant sub-operation in both proving and verification.

\section{Benchmark Results}

\textbf{Base Operations and Proof Size:}

\begin{table}[H]
\centering
\caption{Base operation costs at $n = 1024$ and proof size.}
\label{tab:bp-base}
\begin{tabular}{lrrr}
\toprule
\textbf{Metric} & \textbf{BLS12-381} & \textbf{BN-254} & \textbf{BLS12-377} \\
\midrule
Scalar mul (ms)              & 0.2768 & 0.1229 & 0.3047 \\
MSM at $n=1024$ (ms)         & 28.02  & 16.62  & 28.12  \\
Proof size at $n=1024$ (bytes) & 2216 & 1544   & 2216   \\
\bottomrule
\end{tabular}
\end{table}

At $n = 1024$ there are $\log_2(1024) = 10$ rounds, each
contributing two curve points $(V_L, V_R)$, plus one final point
$\bar{G}$ and one scalar $\bar{u}$. For BLS12-381 this gives
$10 \times 2 \times 96 + 96 + 32 = 2016 + 128 = 2216$ bytes.
For BN-254: $10 \times 2 \times 64 + 64 + 32 = 1280 + 192 = 1544$ bytes.
Each time $n$ doubles, one more round is added, increasing the proof
by $2 \times 96 = 192$ bytes for BLS12-381 or $2 \times 64 = 128$
bytes for BN-254. The full proof size table across all tested sizes
is shown in Table~\ref{tab:bp-proof-size}.

\begin{table}[H]
\centering
\caption{Bulletproofs proof size in bytes across vector sizes.}
\label{tab:bp-proof-size}
\begin{tabular}{rrrr}
\toprule
$n$ & \textbf{BLS12-381} & \textbf{BN-254} & \textbf{BLS12-377} \\
\midrule
4    & 552  & 392  & 552  \\
8    & 760  & 536  & 760  \\
16   & 968  & 680  & 968  \\
32   & 1176 & 824  & 1176 \\
64   & 1384 & 968  & 1384 \\
128  & 1592 & 1112 & 1592 \\
256  & 1800 & 1256 & 1800 \\
512  & 2008 & 1400 & 2008 \\
1024 & 2216 & 1544 & 2216 \\
\bottomrule
\end{tabular}
\end{table}

\textbf{Setup and Commit Time:}

\begin{table}[H]
\centering
\caption{Setup and commit time in milliseconds.}
\label{tab:bp-setup-commit}
\begin{tabular}{rcccccc}
\toprule
& \multicolumn{3}{c}{\textbf{Setup (ms)}} & \multicolumn{3}{c}{\textbf{Commit (ms)}} \\
\cmidrule(lr){2-4} \cmidrule(lr){5-7}
$n$ & \textbf{381} & \textbf{254} & \textbf{377}
    & \textbf{381} & \textbf{254} & \textbf{377} \\
\midrule
4    & 0.65   & 0.09  & 0.62   & 0.76  & 0.40  & 0.63  \\
8    & 1.30   & 0.20  & 1.12   & 0.99  & 0.50  & 0.94  \\
16   & 2.19   & 0.38  & 2.35   & 1.42  & 0.70  & 1.62  \\
32   & 4.25   & 0.64  & 4.51   & 2.30  & 1.46  & 2.51  \\
64   & 8.05   & 1.22  & 9.67   & 3.92  & 1.86  & 4.34  \\
128  & 17.27  & 4.92  & 25.22  & 6.58  & 4.47  & 6.99  \\
256  & 37.49  & 5.46  & 37.06  & 12.98 & 5.25  & 13.68 \\
512  & 79.54  & 8.63  & 74.28  & 21.57 & 8.82  & 23.36 \\
1024 & 154.91 & 18.36 & 149.98 & 39.05 & 15.41 & 39.98 \\
\bottomrule
\end{tabular}
\end{table}

Bulletproofs setup at $n = 1024$ takes only 155 ms on BLS12-381,
compared to over 1 second for KZG setup at the same size. BN-254
setup is exceptionally fast at 18 ms because its random point
generation is cheaper. Commit cost is an MSM of size $n$ and grows
sub-linearly thanks to Pippenger's algorithm.
\\
\textbf{Prove and Verify Time:}

\begin{table}[H]
\centering
\caption{Prove and verify time in milliseconds.}
\label{tab:bp-prove-verify}
\begin{tabular}{rcccccc}
\toprule
& \multicolumn{3}{c}{\textbf{Prove (ms)}} & \multicolumn{3}{c}{\textbf{Verify (ms)}} \\
\cmidrule(lr){2-4} \cmidrule(lr){5-7}
$n$ & \textbf{381} & \textbf{254} & \textbf{377}
    & \textbf{381} & \textbf{254} & \textbf{377} \\
\midrule
4    & 3.34   & 1.70   & 3.01   & 2.53   & 1.46   & 2.73   \\
8    & 5.63   & 3.27   & 6.43   & 4.81   & 2.66   & 5.76   \\
16   & 10.66  & 6.38   & 11.77  & 8.37   & 5.59   & 9.14   \\
32   & 20.48  & 11.92  & 22.07  & 17.85  & 9.15   & 17.42  \\
64   & 38.90  & 21.81  & 43.59  & 29.77  & 17.10  & 34.16  \\
128  & 76.16  & 43.36  & 91.91  & 57.69  & 38.40  & 63.68  \\
256  & 142.95 & 84.38  & 158.03 & 116.46 & 64.82  & 120.30 \\
512  & 302.36 & 162.33 & 295.19 & 254.48 & 133.71 & 239.04 \\
1024 & 650.35 & 415.67 & 582.06 & 496.26 & 329.64 & 499.11 \\
\bottomrule
\end{tabular}
\end{table}

Prove and verify times both roughly double with each doubling of
$n$, confirming $O(n)$ complexity. At $n = 1024$, proving takes
650 ms on BLS12-381 and verifying takes 496 ms, compared to 42 ms
and 3.8 ms respectively for KZG.

\begin{figure}[H]
\centering
\includegraphics[width=\textwidth]{chapter5/combined_bls12_381}
\caption{Bulletproofs benchmark results on BLS12-381: setup, commit,
prove, and verify time (ms) vs.\ vector size $n$ (log--log scale).}
\label{fig:bp-combined-bls12381}
\end{figure}

\section{Discussion and Summary}

\subsection{Fiat-Shamir: Theory vs Implementation}

The original protocol is interactive. The implementation's
non-interactive Fiat-Shamir transformation via SHA-256 is strictly
better for practice. The proof is a self-contained object that
requires no communication channel. The security of this
transformation holds in the random oracle model.

\subsection{MSM Optimisation Impact}

Without Pippenger's MSM, proving at $n = 1024$ would require
$\sum_{k=1}^{10} 2 \times 2^k \approx 4094$ individual scalar
multiplications. With Pippenger's algorithm each MSM of size $m$
costs roughly $O(m / \log m)$ group operations instead of $O(m)$.
At $n = 1024$ this gives a theoretical speedup factor of around
$\log_2(1024)/2 = 5$. The measured MSM time at $n = 1024$ is
28 ms on BLS12-381, while 1024 naive scalar multiplications would
cost roughly $1024 \times 0.277 \approx 284$ ms, confirming a
real speedup of about 10x.

\subsection{Transparency vs Verification Cost}

Bulletproofs are the right choice when a trusted setup is not
acceptable. At $n = 1024$, verifying takes 496 ms on BLS12-381,
about 130x slower than KZG's 3.8 ms which is almost constant. For on-chain verifiers or
high-throughput systems this is prohibitive. For small $n$ or
offline verification the cost is acceptable.

\subsection{Curve Comparison}

BN-254 is the fastest across all metrics. The gap is larger than
in KZG or MKZG because Bulletproofs involve far more scalar
multiplications, which directly track the per-curve scalar
multiplication speed. BLS12-381 and BLS12-377 perform similarly,
with BLS12-377 slightly slower at larger sizes.

\begin{table}[H]
\centering
\caption{Bulletproofs performance summary at $n = 1024$.}
\label{tab:bp-summary}
\begin{tabular}{lrrr}
\toprule
\textbf{Metric} & \textbf{BLS12-381} & \textbf{BN-254} & \textbf{BLS12-377} \\
\midrule
Setup (ms)         & 154.91 & 18.36  & 149.98 \\
Commit (ms)        & 39.05  & 15.41  & 39.98  \\
Prove (ms)         & 650.35 & 415.67 & 582.06 \\
Verify (ms)        & 496.26 & 329.64 & 499.11 \\
Scalar mul (ms)    & 0.2768 & 0.1229 & 0.3047 \\
MSM (ms)           & 28.02  & 16.62  & 28.12  \\
Proof size (bytes) & 2216   & 1544   & 2216   \\
\bottomrule
\end{tabular}
\end{table}